\documentclass[10pt,a4paper]{amsart}
\usepackage{amsmath,amssymb,amsthm,graphicx}
\usepackage{verbatim}
\newenvironment{reference}{\comment}{\endcomment}
%\newenvironment{reference}{}{}
\newenvironment{slogan}{\comment}{\endcomment}
\newenvironment{history}{\comment}{\endcomment}

% For commutative diagrams we use Xy-pic
\usepackage[all]{xy}

% We use 2cell for 2-commutative diagrams.
\xyoption{2cell}
% Theorem environments.
%
\theoremstyle{plain}
\newtheorem{theorem}[subsection]{Theorem}
\newtheorem{proposition}[subsection]{Proposition}
\newtheorem{lemma}[subsection]{Lemma}

\theoremstyle{definition}
\newtheorem{definition}[subsection]{Definition}
\newtheorem{example}[subsection]{Example}
\newtheorem{exercise}[subsection]{Exercise}
\newtheorem{situation}[subsection]{Situation}

\theoremstyle{remark}
\newtheorem{remark}[subsection]{Remark}
\newtheorem{remarks}[subsection]{Remarks}

\numberwithin{equation}{subsection}

% Macros
%
\def\lim{\mathop{\mathrm{lim}}\nolimits}
\def\colim{\mathop{\mathrm{colim}}\nolimits}
\def\Spec{\mathop{\mathrm{Spec}}}
\def\Hom{\mathop{\mathrm{Hom}}\nolimits}
\def\Ext{\mathop{\mathrm{Ext}}\nolimits}
\def\SheafHom{\mathop{\mathcal{H}\!\mathit{om}}\nolimits}
\def\SheafExt{\mathop{\mathcal{E}\!\mathit{xt}}\nolimits}
\def\Sch{\mathit{Sch}}
\def\Mor{\mathop{\mathrm{Mor}}\nolimits}
\def\Ob{\mathop{\mathrm{Ob}}\nolimits}
\def\Sh{\mathop{\mathit{Sh}}\nolimits}
\def\NL{\mathop{N\!L}\nolimits}
\def\CH{\mathop{\mathrm{CH}}\nolimits}
\def\proetale{{pro\text{-}\acute{e}tale}}
\def\etale{{\acute{e}tale}}
\def\QCoh{\mathit{QCoh}}
\def\Ker{\mathop{\mathrm{Ker}}}
\def\Im{\mathop{\mathrm{Im}}}
\def\Coker{\mathop{\mathrm{Coker}}}
\def\Coim{\mathop{\mathrm{Coim}}}

% Boxtimes
%
\DeclareMathSymbol{\boxtimes}{\mathbin}{AMSa}{"02}

%
% Macros for moduli stacks/spaces
%
\def\QCohstack{\mathcal{QC}\!\mathit{oh}}
\def\Cohstack{\mathcal{C}\!\mathit{oh}}
\def\Spacesstack{\mathcal{S}\!\mathit{paces}}
\def\Quotfunctor{\mathrm{Quot}}
\def\Hilbfunctor{\mathrm{Hilb}}
\def\Curvesstack{\mathcal{C}\!\mathit{urves}}
\def\Polarizedstack{\mathcal{P}\!\mathit{olarized}}
\def\Complexesstack{\mathcal{C}\!\mathit{omplexes}}
% \Pic is the operator that assigns to X its picard group, usage \Pic(X)
% \Picardstack_{X/B} denotes the Picard stack of X over B
% \Picardfunctor_{X/B} denotes the Picard functor of X over B
\def\Pic{\mathop{\mathrm{Pic}}\nolimits}
\def\Picardstack{\mathcal{P}\!\mathit{ic}}
\def\Picardfunctor{\mathrm{Pic}}
\def\Deformationcategory{\mathcal{D}\!\mathit{ef}}
\usepackage{cleveref}
\crefname{originallemma}{original lemma}{original lemmas}
\Crefname{originallemma}{Original lemma}{Original lemmas}
\makeatletter
\newlabel{lemma-localized-chern-pre}{{0F9K}{source}{}{sourcepointer-1}{}}
\newlabel{lemma-localized-chern-pre@cref}{{[originallemma][0][]0F9K}{[1][0][]source}{}{}{}}
\makeatother
\crefname{originalreference}{original reference}{original references}
\Crefname{originalreference}{Original reference}{Original references}
\makeatletter
\newlabel{situation-loc-chern}{{0GUJ}{source}{}{sourcepointer-2}{}}
\newlabel{situation-loc-chern@cref}{{[originalreference][0][]0GUJ}{[1][0][]source}{}{}{}}
\makeatother
\crefname{originaldefinition}{original definition}{original definitions}
\Crefname{originaldefinition}{Original definition}{Original definitions}
\makeatletter
\newlabel{definition-localized-chern}{{0FB5}{source}{}{sourcepointer-3}{}}
\newlabel{definition-localized-chern@cref}{{[originaldefinition][0][]0FB5}{[1][0][]source}{}{}{}}
\makeatother
\crefname{originallemma}{original lemma}{original lemmas}
\Crefname{originallemma}{Original lemma}{Original lemmas}
\makeatletter
\newlabel{lemma-gysin-at-infty}{{0F9J}{source}{}{sourcepointer-4}{}}
\newlabel{lemma-gysin-at-infty@cref}{{[originallemma][0][]0F9J}{[1][0][]source}{}{}{}}
\makeatother
\crefname{originallemma}{original lemma}{original lemmas}
\Crefname{originallemma}{Original lemma}{Original lemmas}
\makeatletter
\newlabel{lemma-loc-chern-agree}{{0FB8}{source}{}{sourcepointer-5}{}}
\newlabel{lemma-loc-chern-agree@cref}{{[originallemma][0][]0FB8}{[1][0][]source}{}{}{}}
\makeatother
\crefname{originallemma}{original lemma}{original lemmas}
\Crefname{originallemma}{Original lemma}{Original lemmas}
\makeatletter
\newlabel{lemma-base-change-loc-chern}{{0FB3}{source}{}{sourcepointer-6}{}}
\newlabel{lemma-base-change-loc-chern@cref}{{[originallemma][0][]0FB3}{[1][0][]source}{}{}{}}
\makeatother
\crefname{originalreference}{original reference}{original references}
\Crefname{originalreference}{Original reference}{Original references}
\makeatletter
\newlabel{remark-gysin-on-cycles}{{0B6Z}{source}{}{sourcepointer-7}{}}
\newlabel{remark-gysin-on-cycles@cref}{{[originalreference][0][]0B6Z}{[1][0][]source}{}{}{}}
\makeatother
\crefname{originallemma}{original lemma}{original lemmas}
\Crefname{originallemma}{Original lemma}{Original lemmas}
\makeatletter
\newlabel{lemma-support-cap-effective-Cartier}{{02T9}{source}{}{sourcepointer-8}{}}
\newlabel{lemma-support-cap-effective-Cartier@cref}{{[originallemma][0][]02T9}{[1][0][]source}{}{}{}}
\makeatother
\title{The localized class computed from a blowing up equals the localized Chern class of the perfect object it restricts to}
\author{The Stacks Project Authors}
\date{stacks-project:chow.tex:local-snapshot-20261012; GNU FDL 1.2 or later}
\begin{document}
\maketitle
\noindent\textit{Editorial note.} All delivered mathematical text is the author's own, copied from the pinned Stacks Project chapter file chow.tex; no statement or proof step was rewritten, reordered or removed.\par\medskip
\noindent\textit{Editorial note.} Preamble provenance: the pinned chapter file opens with an input of the shared preamble and carries no preamble text of its own, so the project's single Stacks-wide preamble staged at sources/stacks-chow/preamble.tex is excerpted verbatim as two spans, lines 14-18 (Xy-pic setup) and lines 154-217 (the native theorem declarations and mathematical macros). Nothing was added, removed or redefined; the wrapper class and the amsmath/amssymb/amsthm/graphicx base packages are supplied by the assembler.\par\medskip
\noindent\textit{Editorial note.} The counted statement and proof contain no cross reference, no citation and no included file, so the delivered body is self contained.\par\medskip
\noindent\textit{Editorial note.} Canonical location of the counted proposition: Stacks tag 0FE8 (sources/stacks-chow/pinned-original-source.tex line 11165).\par\medskip
\noindent\textit{Allowed standard prerequisites (editorial).}\begin{itemize}
\item The localized Chern class construction and its description through E and C. Same chapter, Lemma labelled lemma-localized-chern-pre, cited at proof lines 11186-11188 and 11200-11206. Fixes c'\_p(Q) and its expression as a pushforward composed with C.
\item Agreement of the localized class on E with the class computed on W\_infinity. Same chapter, Lemma labelled lemma-loc-chern-agree, cited at proof lines 11208-11211. Performs the second step of the chain.
\item Base change of localized classes and the behaviour of the class C. Same chapter, Lemmas labelled lemma-base-change-loc-chern and lemma-gysin-at-infty, cited at proof lines 11213-11220 and 11190-11194. Supplies the third step and the cycle beta used in place of alpha.
\item Annihilation of cycles supported at infinity and rational equivalence of the special and generic fibres. Same chapter, Remark labelled remark-gysin-on-cycles and Lemma labelled lemma-support-cap-effective-Cartier, cited at proof lines 11230-11234 and 11236-11242. Justifies the fifth and sixth equalities of the chain.
\end{itemize}\medskip
\ifdefined\hypertarget\hypertarget{sourcepointer-1}{}\fi
\noindent\textit{Source reference.} Original lemma 0F9K: sources/stacks-chow/pinned-original-source.tex, source lines 10204--10204. Same-chapter pointer into the author environment lemma at source line 10204; invoked at proof line 11218 of the counted proof. It is retained only as a navigation pointer; every load-bearing step is inside the delivered span.\par\medskip
\ifdefined\hypertarget\hypertarget{sourcepointer-2}{}\fi
\noindent\textit{Source reference.} Original reference 0GUJ: sources/stacks-chow/pinned-original-source.tex, source lines 10678--10678. Same-chapter pointer into the author environment situation at source line 10678; invoked at proof line None of the counted proof. It is retained only as a navigation pointer; every load-bearing step is inside the delivered span.\par\medskip
\ifdefined\hypertarget\hypertarget{sourcepointer-3}{}\fi
\noindent\textit{Source reference.} Original definition 0FB5: sources/stacks-chow/pinned-original-source.tex, source lines 10787--10787. Same-chapter pointer into the author environment definition at source line 10787; invoked at proof line None of the counted proof. It is retained only as a navigation pointer; every load-bearing step is inside the delivered span.\par\medskip
\ifdefined\hypertarget\hypertarget{sourcepointer-4}{}\fi
\noindent\textit{Source reference.} Original lemma 0F9J: sources/stacks-chow/pinned-original-source.tex, source lines 10026--10026. Same-chapter pointer into the author environment lemma at source line 10026; invoked at proof line 11192 of the counted proof. It is retained only as a navigation pointer; every load-bearing step is inside the delivered span.\par\medskip
\ifdefined\hypertarget\hypertarget{sourcepointer-5}{}\fi
\noindent\textit{Source reference.} Original lemma 0FB8: sources/stacks-chow/pinned-original-source.tex, source lines 10975--10975. Same-chapter pointer into the author environment lemma at source line 10975; invoked at proof line 11219 of the counted proof. It is retained only as a navigation pointer; every load-bearing step is inside the delivered span.\par\medskip
\ifdefined\hypertarget\hypertarget{sourcepointer-6}{}\fi
\noindent\textit{Source reference.} Original lemma 0FB3: sources/stacks-chow/pinned-original-source.tex, source lines 10824--10824. Same-chapter pointer into the author environment lemma at source line 10824; invoked at proof line 11222 of the counted proof. It is retained only as a navigation pointer; every load-bearing step is inside the delivered span.\par\medskip
\ifdefined\hypertarget\hypertarget{sourcepointer-7}{}\fi
\noindent\textit{Source reference.} Original reference 0B6Z: sources/stacks-chow/pinned-original-source.tex, source lines 5237--5237. Same-chapter pointer into the author environment remark at source line 5237; invoked at proof line 11233 of the counted proof. It is retained only as a navigation pointer; every load-bearing step is inside the delivered span.\par\medskip
\ifdefined\hypertarget\hypertarget{sourcepointer-8}{}\fi
\noindent\textit{Source reference.} Original lemma 02T9: sources/stacks-chow/pinned-original-source.tex, source lines 5138--5138. Same-chapter pointer into the author environment lemma at source line 5138; invoked at proof line 11240 of the counted proof. It is retained only as a navigation pointer; every load-bearing step is inside the delivered span.\par\medskip
\begin{lemma}
\label{lemma-independent-loc-chern-bQ}
Assume $(S, \delta), X, Z, b : W \to \mathbf{P}^1_X, Q, T, p$
satisfy the assumptions of Lemma \ref{lemma-localized-chern-pre}.
Let $F \in D(\mathcal{O}_X)$ be a perfect object such that
\begin{enumerate}
\item the restriction of $Q$ to $b^{-1}(\mathbf{A}^1_X)$ is
isomorphic to the pullback of $F$,
\item $F|_{X \setminus Z}$ is zero, resp.\ isomorphic to a finite
locally free $\mathcal{O}_{X \setminus Z}$-module of rank $< p$
sitting in cohomological degree $0$, and
\item $Q$ on $W$ and $F$ on $X$ satisfy assumption (3) of
Situation \ref{situation-loc-chern}.
\end{enumerate}
Then the class $P'_p(Q)$, resp.\ $c'_p(Q)$ in $A^p(Z \to X)$ constructed
in Lemma \ref{lemma-localized-chern-pre}
is equal to $P_p(Z \to X, F)$, resp.\ $c_p(Z \to X, F)$
from Definition \ref{definition-localized-chern}.
\end{lemma}

\begin{proof}
The assumptions are preserved by base change with a morphism
$X' \to X$ locally of finite type. Hence it suffices to show that
$P_p(Z \to X, F) \cap \alpha = P'_p(Q) \cap \alpha$,
resp.\ $c_p(Z \to X, F) \cap \alpha = c'_p(Q) \cap \alpha$
for any $\alpha \in \CH_k(X)$. Choose $\beta \in \CH_{k + 1}(W)$
whose restriction to $b^{-1}(\mathbf{A}^1_X)$ is equal to
the flat pullback of $\alpha$ as in the construction of
$C$ in Lemma \ref{lemma-gysin-at-infty}.
Denote $W' = b^{-1}(Z)$ and denote
$E = W'_\infty \subset W_\infty$ the inverse image of $Z$
by $W_\infty \to X$.
The lemma follows from
the following sequence of equalities (the case of $P_p$ is similar)
\begin{align*}
c'_p(Q) \cap \alpha
& =
(E \to Z)_*(c'_p(Q|_E) \cap i_\infty^*\beta) \\
& =
(E \to Z)_*(c_p(E \to W_\infty, Q|_{W_\infty}) \cap i_\infty^*\beta) \\
& =
(W'_\infty \to Z)_*(c_p(W' \to W, Q) \cap i_\infty^*\beta) \\
& =
(W'_\infty \to Z)_*((i'_\infty)^*(c_p(W' \to W, Q) \cap \beta)) \\
& =
(W'_\infty \to Z)_*((i'_\infty)^*(c_p(Z' \to X, F) \cap \beta)) \\
& =
(W'_0 \to Z)_*((i'_0)^*(c_p(Z' \to X, F) \cap \beta)) \\
& =
(W'_0 \to Z)_*(c_p(Z' \to X, F) \cap i_0^*\beta)) \\
& =
c_p(Z \to X, F) \cap \alpha
\end{align*}
The first equality is the construction of $c'_p(Q)$ in
Lemma \ref{lemma-localized-chern-pre}.
The second is Lemma \ref{lemma-loc-chern-agree}.
The base change of $W' \to W$ by $W_\infty \to W$ is the
morphism $E = W'_\infty \to W_\infty$. Hence the third equality holds
by Lemma \ref{lemma-base-change-loc-chern}. The fourth
equality, in which $i'_\infty : W'_\infty \to W'$ is the
inclusion morphism, follows from the fact that $c_p(W' \to W, Q)$
is a bivariant class. For the fifth equality, observe that
$c_p(W' \to W, Q)$ and $c_p(Z' \to X, F)$
restrict to the same bivariant class in
$A^p((b')^{-1} \to b^{-1}(\mathbf{A}^1_X))$ by
assumption (1) of the lemma which says that $Q$ and $F$ restrict
to the same object of $D(\mathcal{O}_{b^{-1}(\mathbf{A}^1_X)})$;
use Lemma \ref{lemma-base-change-loc-chern}.
Since $(i'_\infty)^*$ annihilates cycles supported on $W'_\infty$
(see Remark \ref{remark-gysin-on-cycles}) we conclude the fifth equality
is true. The sixth equality holds because $W'_\infty$ and $W'_0$
are the pullbacks of the rationally equivalent effective Cartier divisors
$D_0, D_\infty$ in $\mathbf{P}^1_Z$ and hence $i_\infty^*\beta$ and
$i_0^*\beta$ map to the same cycle class on $W'$; namely, both
represent the class
$c_1(\mathcal{O}_{\mathbf{P}^1_Z}(1)) \cap c_p(Z \to X, F_) \cap \beta$ by
Lemma \ref{lemma-support-cap-effective-Cartier}.
The seventh equality holds because $c_p(Z \to X, F)$ is
a bivariant class. By construction $W'_0 = Z$ and $i_0^*\beta = \alpha$
which explains why the final equality holds.
\end{proof}

\end{document}
